% ============================================================
% Reporte: mapeo-memoria
% Generado a partir de reportes/comun/plantilla_base.tex
% (NO edites la plantilla: cada unidad trabaja en su main.tex)
%
% Reglas de secciones y extensiones: reportes/AGENTS.md
% Compilar SIEMPRE con:
%   ../../scripts/latexbuild main.tex      (desde esta carpeta)
%   → genera build/main.pdf
% ============================================================
\documentclass[llncs]{llncs}

\usepackage[utf8]{inputenc}
\usepackage[T1]{fontenc}

% ---- AVISO DE CODIFICACIÓN (xelatex + T1) ----
% Con la cadena del proyecto (xelatex con codificación T1) los caracteres
% UTF-8 de la banda 0x80-0xFF escritos como literal se remapean mal y
% salen otros glifos: « » § ª · aparecen como ń ż ğ ł ů.
% En su lugar usa las macros de LaTeX, que sí funcionan:
%   «texto»  ->  ``texto''          §6.2   ->  \S6.2
%   ª        ->  \textordfeminine   ·      ->  \textperiodcentered
% Sí funcionan como literal: los acentos (á é í ó ú ñ ü) y  – — “ ” → .
\usepackage{graphicx}
\usepackage{listings}
\usepackage{hyperref}
\usepackage{xcolor}

% ---- Estilo para código VHDL (apéndice y fragmentos) ----
\lstdefinestyle{vhdl}{
  language=VHDL,
  basicstyle=\ttfamily\scriptsize,
  keywordstyle=\color{black},
  commentstyle=\color{gray}\itshape,
  stringstyle=\color{black},
  showstringspaces=false,
  upquote=true,
  columns=fullflexible,
  breaklines=true,
  frame=single,
  rulecolor=\color{black}
}
\lstset{style=vhdl}
\renewcommand{\lstlistingname}{Listado}

% ---- Estilo para bloques de comandos del repositorio (sección 3) ----
\lstdefinestyle{shell}{
  basicstyle=\ttfamily\scriptsize,
  keywordstyle=\color{black},
  commentstyle=\color{gray}\itshape,
  stringstyle=\color{black},
  showstringspaces=false,
  upquote=true,
  columns=fullflexible,
  breaklines=true,
  frame=single,
  rulecolor=\color{black}
}

\begin{document}

% ============================================================
% TÍTULO Y AUTORES  (los 5 integrantes del equipo — por rellenar)
% ============================================================
\title{Comparación de las funciones de mapeo de la memoria caché\\
mediante simulación VHDL}

\author{Alfredo Ramírez Candia$^{1}$%
\and Aldo Mauricio Villegas Peralta$^{1}$%
\and Uriel Lemus Pérez$^{1}$%
\and Santiago Silverio Flores$^{1}$%
\and Max Ramos Martinez$^{1}$}

\authorrunning{Apellido et al.}
\institute{Equipo de Arquitectura de Computadoras,\\
Benemérita Universidad Autónoma de Puebla, Puebla, México.\\
\email{\{rc202375947, vp202368617, lp202362797, sf202367190, rm202376182\}@alm.buap.mx}}

\maketitle

% ============================================================
% RESUMEN  (< 100 palabras — fuente 9pt con sangría, estilo llncs)
% ============================================================
\begin{abstract}
Resumen. Se comparan mediante simulación VHDL-2008 las tres funciones
de mapeo de la memoria caché: directo, asociativo total y asociativo
por conjuntos. Las tres cachés comparten interfaz, memoria principal
y una traza de 20 accesos con un patrón de thrashing, de modo que la
comparación es justa. Las simulaciones se ejecutan con GHDL y las
ondas se inspeccionan con GTKWave. En la traza común la caché directa
alcanza 35\,\% de aciertos, la por conjuntos 50\,\% y la asociativa
total 65\,\%, resultados coherentes con el coste y la complejidad de
cada mapeo.
\end{abstract}

\noindent\textbf{Palabras Clave:} Memoria caché, Mapeo directo,
Asociativo por conjuntos, VHDL, GHDL, Simulación

% ============================================================
% 1. INTRODUCCIÓN  (1–2 cuartillas)
% ============================================================
\section{Introducción}
\label{sec:intro}

La memoria principal es mucho más lenta que el procesador: cerrar una
brecha de este tipo es el motivo de existir de la jerarquía de
memoria, y la caché es su primer y más efectivo nivel
\cite{stallings2019,hennessy2019,tanenbaum2013}. Una caché aprovecha
los principios de localidad temporal y espacial \cite{wikipedia-locality}
para guardar en un almacenamiento
rápido y pequeño una copia de los bloques que el procesador usa con
más frecuencia \cite{patterson2017,bryant2015}.

Cuando se diseña una caché hay que decidir, para cada bloque de la
memoria principal, qué línea de la caché puede alojarlo. Esa decisión
es la \emph{función de mapeo}, y existen tres básicos: mapeo directo
(\emph{direct mapping}), mapeo asociativo total (\emph{associative
mapping}) y mapeo asociativo por conjuntos (\emph{set-associative
mapping})
\cite{stallings2019,wikipedia-cache,null2024}. Cada uno compromete
distinto hardware de comparación y distinta flexibilidad de
reemplazo, de modo que ninguna opción es \emph{a priori} mejor que
las otras: dependen del patrón de accesos.

Este trabajo construye y compara las tres funciones de mapeo en el
mismo entorno de simulación. Las tres cachés están descritas en
VHDL-2008~\cite{ieee1076}, se simulan con GHDL~\cite{ghdl-docs} y sus ondas
se
inspeccionan con GTKWave~\cite{gtkwave-docs}. Para que la comparación
sea justa, los tres módulos exponen \emph{la misma interfaz}, leen de
\emph{la misma memoria principal de comportamiento} y ejecutan
\emph{la misma traza de accesos}, que incluye un patrón de
\emph{thrashing} típico del mapeo directo. El objetivo es doble:
(i) comprobar por simulación el comportamiento teórico de cada
mapeo, y (ii) cuantificar sus diferencias en aciertos y fallos sobre
una carga común.

Desde el punto de vista didáctico conviene aclarar la terminología. El
equipo denomina a los tres métodos ``directo'', ``asociativo'' y
``asociativo por conjuntos'', siguiendo la descripción de funcionamiento
del laboratorio (búsqueda lineal de instrucciones, uso de etiquetas
con comparación en paralelo y división en conjuntos); en la
bibliografía se usa la terminología estándar de Stallings, que es la
que se emplea en el resto del artículo \cite{stallings2019}.

El artículo se estructura como sigue. La sección~\ref{sec:diseno}
describe las tres funciones de mapeo y los módulos VHDL. La
sección~\ref{sec:propuesta} presenta la metodología experimental. La
sección~\ref{sec:simulaciones} muestra las capturas de las
simulaciones, la sección~\ref{sec:resultados} analiza los resultados
cuantitativos y la sección~\ref{sec:conclusiones} concluye y señala
trabajos futuros.

% ============================================================
% 2. DISEÑO TEÓRICO  (3–5 cuartillas)
% ============================================================
\section{Diseño Teórico}
\label{sec:diseno}

\subsection{Contexto: memoria caché y bloques}

Una caché organiza la memoria principal en \emph{bloques} (o
líneas) de varios bytes y guarda esos bloques en sus \emph{líneas}.
En este trabajo se reproducen a escala reducida los parámetros del
Ejemplo 5.1 del capítulo 5 de Stallings \cite{stallings2019} (caché de
64\,kB, memoria principal de 16\,MB y bloques de 4 bytes):
direccionamiento a nivel de byte
con direcciones de 10 bits, palabra y bloque de 4 bytes (por tanto 2
bits de offset), memoria principal de 1\,KiB (256 bloques) y caché de
64\,B (16 líneas). La Tabla~\ref{tab:par} recoge los parámetros y el
reparto de bits de la dirección de cada mapeo.

\begin{table}[t]
\centering
\caption{Parámetros comunes y reparto de la dirección (10 bits).}
\label{tab:par}
\small
\begin{tabular}{llll}
\hline\noalign{\smallskip}
Parámetro & Valor & Mapeo & Reparto de bits\\
\noalign{\smallskip}\hline\noalign{\smallskip}
Dirección & 10 bits (byte) & Directo & tag(4) $\vert$ índice(4) $\vert$ offset(2)\\
Palabra/bloque & 4 bytes & Asociativo & tag(8) $\vert$ offset(2)\\
Mem. principal & 1\,KiB (256 bloques) & Conjuntos & tag(5) $\vert$ conjunto(3) $\vert$ offset(2)\\
Caché & 64\,B (16 líneas) & & \\
\noalign{\smallskip}\hline
\end{tabular}
\end{table}

El papel de la caché se entiende dentro de la jerarquía de memoria:
cuanto más cerca del procesador, los almacenamientos son más rápidos y
más pequeños, y más lentos y grandes al alejarse
\cite{tanenbaum2013,hamacher2002,wikipedia-locality}. La caché es el primer nivel
de esa jerarquía y vive de los principios de localidad temporal y
espacial \cite{wikipedia-locality,stallings2019,smith1982}: un programa vuelve
a usar pronto los mismos datos y usa datos que guardan relación entre
sí. Para que la localidad espacial rinda, el bloque debería contener
varias palabras consecutivas; aquí el bloque es de una sola palabra
(4 bytes), una simplificación deliberada que aísla el efecto del
mapeo y que se discute en la sección~\ref{sec:resultados}.

\begin{figure}[t]
  \centering
  \includegraphics[width=0.85\textwidth]{figuras/reparto_direccion.png}
  \caption{Reparto de la dirección de 10 bits en los tres mapeos: la
    etiqueta se guarda, el índice o el conjunto sitúa la línea y el
    offset elige el byte dentro del bloque.}
  \label{fig:reparto}
\end{figure}

La Fig.~\ref{fig:reparto} ilustra el reparto de los bits: al crecer
la asociatividad, los bits de índice pasan a formar parte de la
etiqueta, y por eso hay que comparar más etiquetas en cada acceso.

\subsection{Mapeo directo}

En el mapeo directo cada bloque solo puede ocupar una única línea de
la caché: la línea $i$ del bloque se calcula como $i = j \bmod m$,
donde $j$ es el número de bloque en memoria y $m$ el número de líneas
de la caché \cite{stallings2019,wikipedia-cpu-cache}. En nuestro caso $m=16$, así
que el índice son los 4 bits intermedios de la dirección y los 4 bits
superiores quedan como etiqueta (\emph{tag}) que debe guardarse y
compararse. La búsqueda consiste en un único comparador de 4 bits
sobre la línea seleccionada por el índice, por lo que es el mapeo más
barato y rápido: no hay que recorrer la caché ni decidir víctima.

Su inconveniente es estructural: dos bloques con el mismo índice
interfieren siempre, y si un programa alterna accesos a dos bloques
que comparten índice, cada acceso es un fallo: es el fenómeno de
\emph{thrashing}
\cite{stallings2019,wikipedia-cache,wikipedia-metrics}. La traza de
este trabajo está diseñada para provocarlo con las palabras 0, 16 y
32, que comparten índice.

\subsection{Mapeo asociativo total}

En el mapeo asociativo total cualquier bloque puede ocupar
cualquiera de las líneas \cite{stallings2019}. La dirección se
divide solo en etiqueta (8 bits) y offset, y la caché guarda la
etiqueta de cada línea: en el ciclo de comparación se comparan
\emph{en paralelo} las 16 etiquetas con la solicitada. Esto requiere
16 comparadores de 8 bits y una red de selección (la estructura
clásica es una memoria de contenido o \emph{CAM}, \emph{content
addressable memory}) \cite{stallings2019,hennessy2019}. El
inconveniente es el coste y la potencia de los comparadores en
paralelo; la ventaja es que un bloque nunca expulsa a otro por
colisión de índice: solo sale de la caché cuando hace falta espacio y
se elige una víctima con una política de reemplazo, aquí LRU
(\emph{least recently used}) \cite{stallings2019}.

\subsection{Mapeo asociativo por conjuntos}

El mapeo por conjuntos es el punto intermedio \cite{stallings2019}.
La caché se divide en conjuntos de $k$ líneas (\emph{$k$-way
set-associative}), y un bloque solo puede ocupar cualquiera de las
$k$ líneas de \emph{su} conjunto; dentro del conjunto la elección es
libre. Con $k=1$ se obtiene el mapeo directo y con $k=m$ el asociativo
total \cite{stallings2019,patterson2017}. En nuestro caso hay 8
conjuntos de 2 vías, de modo que la dirección se reparte en etiqueta
(5 bits), conjunto (3 bits) y offset, y solo hacen falta 2
comparadores de 5 bits por acceso. El reemplazo es LRU con 1 bit por
conjunto \cite{wikipedia-cpu-cache}: se recuerda cuál de las dos vías se usó
más recientemente y la víctima es la otra.

\subsection{Reemplazo LRU}

Las cachés asociativas necesitan una política de reemplazo. Se ha
elegido LRU, que descarta el bloque hace más tiempo sin usar y que
aprovecha mejor la localidad temporal que políticas aleatorias o FIFO
\cite{stallings2019,hennessy2019,wikipedia-replacement}. La caché totalmente
asociativa
implementa LRU exacto con un contador de antigüedad por línea (0 =
más reciente), y la por conjuntos implementa LRU exacto con 1 bit por
conjunto al tener solo dos vías. En un diseño real el LRU exacto de
16 vías no es barato y se sustituye por aproximaciones
\cite{hennessy2019}; aquí se justifica por tratarse de un modelo de
simulación de tamaño reducido. Las alternativas a LRU son bien
conocidas: aleatorio, FIFO o aproximaciones como pseudo-LRU, que
apenas pierden tasa de aciertos y ahorran coste
\cite{hennessy2019,wikipedia-replacement}. Aquí solo hacen falta dos
variantes baratas: un contador de antigüedad por línea en la
totalmente asociativa y un bit por conjunto en la por conjuntos, que
con dos vías equivale a LRU exacto \cite{wikipedia-cpu-cache}. Conviene
recordar, además, por qué no se usa FIFO: su sustituto óptimo puede
empeorar al añadir marcos, la conocida anomalía de Belady
\cite{belady1969}.

\subsection{Funcionamiento de una consulta}

Los tres cachés comparten interfaz y máquina de estados, de modo que
una consulta recorre los mismos pasos en los tres; solo cambia cómo
se busca y se compara la etiqueta.

\textbf{Acierto.} El banco de prueba pone la dirección en
\texttt{addr} y activa \texttt{req}. En el flanco siguiente la
máquina sale de \texttt{IDLE} y entra en \texttt{COMPARE}: se extraen
los campos de la dirección y se comparan la etiqueta (o las
etiquetas) con la de la línea seleccionada; con coincidencia,
\texttt{hit} se activa y el dato de la línea queda disponible. El
flanco siguiente la máquina pasa a \texttt{DONE}, donde se levanta
\texttt{ready}; el banco toma \texttt{data} y baja \texttt{req}. Un
acierto ocupa tres ciclos de reloj.

\textbf{Fallo.} El comienzo es idéntico, pero sin coincidencia
\texttt{hit} permanece baja y la máquina avanza a
\texttt{FETCH\_MEM}: la caché emite la dirección del bloque por
\texttt{mem\_addr} y activa \texttt{mem\_req}. La memoria principal,
de latencia de un ciclo, responde con \texttt{mem\_ready} y el dato;
entonces la caché guarda el bloque en su línea, la marca como válida,
actualiza su política de reemplazo y pasa a \texttt{DONE}, donde
\texttt{ready} entrega el dato solicitado. El paso por
\texttt{FETCH\_MEM} y la espera a la memoria alargan el fallo hasta
seis ciclos. La sección~\ref{sec:resultados} mide ambas latencias
sobre los ficheros~VCD.

\subsection{Coste de cada mapeo}

El coste de un mapeo se concentra en comparar etiquetas y en decidir
víctimas. La Tabla~\ref{tab:coste} resume lo que implica cada uno
para la caché de este trabajo.

\begin{table}[t]
\centering
\small
\begin{tabular}{@{}lrrr@{}}
\hline\noalign{\smallskip}
Mapeo & Comparadores & Reemplazo & Selección \\
\noalign{\smallskip}\hline\noalign{\smallskip}
Directo & 1 $\times$ 4 bits & ninguno & por índice \\
Asociativo total & 16 $\times$ 8 bits & contador de 4 bits & mux 16 a 1 \\
Por conjuntos & 2 $\times$ 5 bits & 1 bit por conjunto & mux 2 a 1 \\
\noalign{\smallskip}\hline
\end{tabular}
\caption{Coste de la comparación de etiquetas y del reemplazo en cada mapeo.}
\label{tab:coste}
\end{table}

El directo no tiene red de selección ni política que decidir; el
total concentra su coste en los 16 comparadores y en buscar la
víctima entre las 16 líneas; y el por conjuntos obtiene la mayor
parte del efecto del total con dos comparadores por acceso
\cite{hennessy2019,wikipedia-cpu-cache}. Cada salto de asociatividad
compra aciertos con hardware de comparación:
Tabla~\ref{tab:resultados} muestra lo que cuesta cada paso en la
traza de este trabajo.

\subsection{Decisiones de diseño y módulos VHDL}

Las decisiones de diseño, todas ellas dirigidas a que la comparación
sea justa, son:
\begin{itemize}
  \item \textbf{Misma interfaz.} Los tres cachés exponen los mismos
    puertos (\texttt{clk}, \texttt{rst}, \texttt{addr}, \texttt{req},
    \texttt{ready}, \texttt{data}, \texttt{hit}, \texttt{estado} y la
    interface de memoria \texttt{mem\_addr}, \texttt{mem\_req},
    \texttt{mem\_data}, \texttt{mem\_ready}), de modo que el banco de
    pruebas y la memoria son intercambiables.
  \item \textbf{Solo lectura.} Se modela el \emph{fetch} de
    instrucciones: no hay bit sucio (\emph{dirty}) ni política de
    escritura \cite{wikipedia-cache-computing}, que queda como trabajo
    futuro (sección~\ref{sec:conclusiones}).
  \item \textbf{Máquina de estados común.}
    $\texttt{IDLE} \rightarrow \texttt{COMPARE} \rightarrow
    (\texttt{FETCH\_MEM} \rightarrow \texttt{DONE}) \mid
    \texttt{DONE}$: se compara, y solo en caso de fallo se pide el
    bloque a la memoria principal, que tiene una latencia de 1 ciclo
    (modelo de comportamiento).
  \item \textbf{Misma memoria principal.} RAM de comportamiento de
    1\,KiB inicializable con un patrón determinista, compartida por
    los bancos de prueba.
\end{itemize}

Los módulos de \texttt{vhdl/mapeo-memoria/src/} son
\texttt{memoria\_principal} (RAM de 256 palabras de 32 bits con
latencia de 1 ciclo y puerto de escritura para los bancos de prueba),
\texttt{cache\_directo} (1 comparador de 4 bits, sin reemplazo),
\texttt{cache\_asociativa} (16 comparadores de 8 bits en paralelo,
contadores LRU) y \texttt{cache\_conjuntos} (2 comparadores de 5
bits, LRU de 1 bit por conjunto). El código completo está en el
apéndice.

% ============================================================
% 3. PROPUESTA EXPERIMENTAL  (1 hoja)
% ============================================================
\section{Propuesta Experimental}
\label{sec:propuesta}

\subsection{Entorno de simulación}

Las descripciones están escritas en VHDL-2008~\cite{ieee1076}, con el
paquete estándar \texttt{std\_logic\_1164}~\cite{ieee1164} que define el
tipo lógico multivalor de todos los puertos y señales, y siguen la
metodología de modelado y bancos de prueba de la bibliografía de
referencia~\cite{harris2008}. Se verifican por simulación funcional con GHDL~\cite{ghdl-docs}, un
simulador libre y multiplataforma que se emplea en tres
fases~\cite{ghdl-invoking}: \emph{análisis} (\texttt{-a}) del código
fuente, \emph{elaboración} (\texttt{-e}) del diseño y \emph{ejecución}
(\texttt{-r}) del banco de prueba. Durante la ejecución puede volcarse
el valor de todas las señales en un fichero~VCD~\cite{ghdl-simulation},
que después se inspecciona con GTKWave~\cite{gtkwave-docs}. En este
trabajo la cadena se ejecuta siempre mediante los \emph{wrappers} del
repositorio, que fijan las opciones de compilación (VHDL-2008,
biblioteca de trabajo y rutas de salida):

\begin{lstlisting}[style=shell]
scripts/vhdl-analizar mapeo-memoria
scripts/vhdl-sim     mapeo-memoria tb_cache_directo
scripts/vhdl-wave    vhdl/mapeo-memoria/salidas/tb_cache_directo.vcd
\end{lstlisting}

\subsection{Bancos de prueba}

La unidad se verifica con cinco bancos de prueba, todos con
\emph{autochequeo}: cada comprobación es un \texttt{assert} que, si no
se cumple, detiene la simulación con \texttt{severity failure} y
describe el problema, y al finalizar cada banco se emite un
\texttt{report} de éxito del tipo
\texttt{"tb\_cache\_directo:\ OK"}. GHDL devuelve un código de salida
distinto de cero cuando un aserto falla, de modo que la integración
continua del repositorio (que analiza y simula los cinco bancos en
cada \emph{pull request}) detecta cualquier regresión.
Tabla~\ref{tab:bancos} resume qué comprueba cada uno.

\begin{table}[t]
\centering
\small
\begin{tabular}{@{}ll@{}}
\hline
\noalign{\smallskip}
Banco de prueba & Qué verifica \\
\noalign{\smallskip}\hline\noalign{\smallskip}
\texttt{tb\_memoria\_principal} & patrón de fábrica, puerto de \\
 & escritura y palabras vecinas \\
\texttt{tb\_cache\_directo} & dato y acierto/fallo por acceso; \\
 & totales 7 aciertos / 13 fallos \\
\texttt{tb\_cache\_asociativa} & ídem; totales 13 aciertos / 7 fallos \\
\texttt{tb\_cache\_conjuntos} & ídem; totales 10 aciertos / 10 fallos \\
\texttt{tb\_comparacion\_mapeos} & los tres mapeos sobre una \\
 & memoria compartida; tabla comparativa \\
\noalign{\smallskip}\hline
\end{tabular}
\caption{Bancos de prueba de la unidad \texttt{mapeo-memoria}.}
\label{tab:bancos}
\end{table}

Los tres bancos de caché ejecutan la misma traza, y el cuarto
banco (\texttt{tb\_\allowbreak comparacion\_\allowbreak mapeos})
instancia las tres cachés sobre una
\emph{única} memoria de comportamiento, compartida mediante un
arbitraje \emph{round-robin} registrado (una concesión por ciclo), de
modo que las diferencias observadas se deben solo al mapeo y no al
modelo de memoria.

\subsection{Traza común y métricas}

La traza tiene 20 accesos y combina localidad temporal y espacial con
un patrón deliberado de \emph{thrashing}: las palabras 0, 16 y 32
(direcciones byte 0, 64 y 128) compiten por la misma línea en el mapeo
directo y por el mismo conjunto en el mapeo por conjuntos. Cada acceso
se contrasta con el resultado teórico: el dato devuelto debe coincidir
con el patrón de fábrica de la memoria (palabra $i = 3i+7$) y el
acierto o fallo debe coincidir con el patrón esperado del banco. La
métrica de referencia es la tasa de aciertos $h$, calculada como el
número de aciertos dividido entre 20, y, por complemento, la tasa de
fallos $1-h$.

Conviene acotar el alcance de la comparación: es \emph{funcional}, en
el sentido de que compara aciertos y fallos sobre la misma traza y el
mismo modelo de memoria, de latencia de un ciclo. No se mide el
tiempo de ejecución en un procesador real y no se contempla síntesis:
GHDL se usa como simulador, y la verificación sigue el flujo habitual
del proyecto (análisis, elaboración, ejecución y revisión de ondas).

% ============================================================
% 4. SIMULACIONES DE VHDL  (capturas de los módulos más importantes)
% ============================================================
\section{Simulaciones de VHDL}
\label{sec:simulaciones}

Los cinco bancos terminan con su mensaje de \texttt{OK} sin disparar
ninguna aserción (Tabla~\ref{tab:bancos}). Las figuras son capturas de
GTKWave sobre los ficheros~VCD de \texttt{vhdl/mapeo-memoria/salidas/},
con la traza completa a la vista (0~ns hasta el final de la simulación).
Cada captura expone las filas del banco: \texttt{clk}, \texttt{rst},
el contador de \texttt{ciclo}, el estímulo (\texttt{addr}, \texttt{req}),
la respuesta (\texttt{ready}, \texttt{hit}, \texttt{data}, \texttt{estado})
y los contadores de aciertos y fallos
(\texttt{caches\_\allowbreak hits}, \texttt{caches\_\allowbreak misses}),
con los vectores en hexadecimal y los contadores en decimal.

\begin{figure}[!t]
  \centering
  \includegraphics[width=0.82\textwidth]{figuras/captura_memoria_principal.png}
  \caption{Memoria principal (\texttt{tb\_memoria\_principal}): patrón
    de fábrica y dos escrituras verificadas.}
  \label{fig:cap-memoria}
\end{figure}

\textbf{Memoria principal.} La Fig.~\ref{fig:cap-memoria} muestra las
cuatro lecturas iniciales del patrón de fábrica (palabras 0, 1, 16 y
255) y las dos escrituras: \texttt{A5A5F00F} sobre la dirección 20 y
\texttt{000000AA} sobre la 64. Tras cada escritura \texttt{wr\_en}
pulsa un ciclo, el banco relee esa dirección y después una palabra
vecina (0 y 4), que debe quedar intacta.

\begin{figure}[!t]
  \centering
  \includegraphics[width=0.82\textwidth]{figuras/captura_cache_directo.png}
  \caption{Mapeo directo (\texttt{tb\_cache\_directo}): traza completa,
    FSM y contadores de aciertos y fallos.}
  \label{fig:cap-directo}
\end{figure}

\textbf{Mapeo directo.} En la Fig.~\ref{fig:cap-directo} las
direcciones 0, 64 y 128 (palabras 0, 16 y 32) compiten por la línea de
índice 0 y \texttt{hit} queda bajo cuando ésta no contiene la palabra
buscada: a lo largo de la traza \texttt{hit} solo pulsa 7 veces. La fila
\texttt{estado} recorre la FSM y en cada fallo se interna en
\texttt{FETCH\_MEM} (código 2), donde \texttt{mem\_req} pide el bloque a
la memoria. Los contadores escalonados acaban en 7 aciertos
(\texttt{0007} en hexadecimal) y 13 fallos (\texttt{000D}).

\begin{figure}[!t]
  \centering
  \includegraphics[width=0.82\textwidth]{figuras/captura_cache_asociativa.png}
  \caption{Asociativa total (\texttt{tb\_cache\_asociativa}): 7 fallos
    obligatorios y 13 aciertos.}
  \label{fig:cap-asociativa}
\end{figure}

\textbf{Mapeo asociativo total.} La Fig.~\ref{fig:cap-asociativa}
repite la traza sobre las 16 líneas totalmente asociativas: solo
aparecen los 7 fallos obligatorios de la primera visita a cada
dirección. En la fila \texttt{estado} esos fallos iniciales se
internan en \texttt{FETCH\_MEM}, pero a partir del acceso 5 la FSM
salta directo de \texttt{COMPARE} a \texttt{DONE}; el contador
\texttt{caches\_\allowbreak misses} se congela en 7 (\texttt{0007}) y
\texttt{caches\_\allowbreak hits} sube hasta 13 (\texttt{000D}). El
patrón de \emph{thrashing} del mapeo directo no aparece aquí: al no
haber conflicto de índice, las palabras 0, 16 y 32 conviven en
líneas distintas.

\begin{figure}[!t]
  \centering
  \includegraphics[width=0.82\textwidth]{figuras/captura_cache_conjuntos.png}
  \caption{Por conjuntos (\texttt{tb\_cache\_conjuntos}): 8 conjuntos
    de 2 vías y fallos de conflicto en el conjunto 0.}
  \label{fig:cap-conjuntos}
\end{figure}

\textbf{Mapeo asociativo por conjuntos.} La
Fig.~\ref{fig:cap-conjuntos} usa 8 conjuntos de 2 vías: a los 7 fallos
obligatorios se suman 3 de conflicto (accesos 17--19), todos en el
conjunto 0, donde compiten las palabras 0, 16 y 32; con solo dos vías
la expulsión la decide la política LRU descrita en la
sección~\ref{sec:diseno}. La fila \texttt{estado} distingue ambos
casos y, en cada fallo, la interfaz con la memoria se hace visible:
\texttt{mem\_req} pulsa una vez y \texttt{mem\_addr} presenta el bloque
alineado (\texttt{044}, \texttt{080} o \texttt{00C} en los fallos de
conflicto). Totales: 10 aciertos y 10 fallos
(\texttt{000A}/\texttt{000A} en los contadores).

\begin{figure}[!t]
  \centering
  \includegraphics[width=0.82\textwidth]{figuras/captura_comparacion_mapeos.png}
  \caption{Los tres mapeos sobre una memoria compartida: aciertos,
    fallos e interfaz de memoria.}
  \label{fig:cap-comparacion}
\end{figure}

\textbf{Comparación en paralelo.} Por último,
Fig.~\ref{fig:cap-comparacion}, el banco
\texttt{tb\_comparacion\_\allowbreak mapeos} superpone sobre una única
memoria compartida las tres señales de acierto y los seis contadores,
agrupados por mapeo: \texttt{hit\_asociativa} frecuente,
\texttt{hit\_conjuntos} intermedia y \texttt{hit\_directo} escasa.
Debajo aparece la memoria compartida (\texttt{mem\_addr},
\texttt{mem\_req}, \texttt{mem\_ready} y \texttt{mem\_data}), atendida
por el arbitraje \emph{round-robin} descrito en la
sección~\ref{sec:propuesta} (las filas \texttt{grant\_s} y
\texttt{prev\_s} registran la concesión del árbitro). Los totales,
7/13, 10/10 y 13/7, coinciden con la Tabla~\ref{tab:bancos}.

% ============================================================
% 5. RESULTADOS  (3–5 cuartillas)
% ============================================================
\section{Resultados}
\label{sec:resultados}

Esta sección analiza las métricas que los cinco bancos verifican en cada
ejecución. Los recuentos de aciertos y fallos proceden de la salida de
\texttt{tb\_comparacion\_\allowbreak mapeos}, que ejecuta los tres
mapeos sobre la
misma traza y una misma memoria de comportamiento, y las latencias de
acceso se miden sobre los ficheros~VCD de la
sección~\ref{sec:simulaciones}. Se usa la terminología habitual: ratio
de aciertos $h$ (aciertos divididos entre accesos) y tasa de fallos
$1-h$ \cite{tanenbaum2013,null2024}.

\subsection{Métricas globales de la traza}

\begin{table}[t]
\centering
\small
\begin{tabular}{@{}lrrrr@{}}
\hline
\noalign{\smallskip}
Métrica & Directo & Por conjuntos & Total \\ 
\noalign{\smallskip}\hline\noalign{\smallskip}
Aciertos (de 20) & 7 & 10 & 13 \\
Fallos & 13 & 10 & 7 \\
Tasa de aciertos & 35\,\% & 50\,\% & 65\,\% \\
Fallos obligatorios & 7 & 7 & 7 \\
Fallos de conflicto & 6 & 3 & 0 \\
Fin de la simulación (ns) & 1415 & 1325 & 1235 \\
\noalign{\smallskip}\hline
\end{tabular}
\caption{Resultados sobre la traza común de 20 accesos (salida de
\texttt{tb\_comparacion\_mapeos} y fin de simulación de cada banco).}
\label{tab:resultados}
\end{table}

\begin{figure}[t]
  \centering
  \includegraphics[width=0.78\textwidth]{figuras/grafico_aciertos_fallos.png}
  \caption{Composición de los 20 accesos de la traza en cada mapeo:
    aciertos y fallos con su tasa.}
  \label{fig:tasas}
\end{figure}

Sobre los mismos 20 accesos, el mapeo directo resuelve 7 aciertos y 13
fallos (35\,\%), el por conjuntos 10 y 10 (50\,\%) y el total 13 y 7
(65\,\%), como recogen la Tabla~\ref{tab:resultados} y la
Fig.~\ref{fig:tasas}. El orden coincide con el que predicen los tres
esquemas de mapeo \cite{stallings2019,null2024}: al pasar de una línea
fija por bloque a dos vías por conjunto y de ahí a 16 líneas libres,
los accesos que colisionaban dejan de hacerlo. La comparación es justa
porque los tres mapeos comparten la misma memoria y el mismo arbitraje
\emph{round-robin} (sección~\ref{sec:propuesta}), de modo que las
diferencias proceden solo del reparto de direcciones \cite{patterson2017}.

Los números no dependen de la inspección manual de las ondas: cada dato
leído y cada resultado acierto/fallo se contrasta con un \texttt{assert}
dentro de la simulación, y los totales de la tabla se verifican al
final de la ejecución (sección~\ref{sec:propuesta}), de modo que un
cambio de diseño que altere un solo acierto detiene la simulación.

\subsection{Clasificación de los fallos}
\label{sub:3c}

Los tres mapeos comparten los mismos 7 fallos obligatorios: son las
primeras visitas a los siete bloques distintos que toca la traza (las
palabras 0, 1, 2, 3, 16, 17 y 32) y ninguna caché puede resolverlas sin
cargar antes el bloque \cite{wikipedia-metrics,wikipedia-cache}. Tampoco
hay fallos de capacidad, porque el conjunto de trabajo (7 bloques de
4~B) cabe con holgura en las 16 líneas de 4~B de la caché
\cite{wikipedia-metrics,denning1968}. Por tanto, la diferencia observada
--6, 3 y 0-- son fallos de conflicto: los que aparecen cuando un bloque
legítimo no encuentra línea libre en su índice o en su conjunto
\cite{wikipedia-metrics,stallings2019}. La Tabla~\ref{tab:resultados}
desglosa los tres tipos, y su reparto explica por sí solo la tabla de
acierto: los tres tipos suman siempre los 13, 10 y 7 fallos medidos.

\subsection{Thrashing bajo control}

En el mapeo directo cada bloque solo puede ocupar su línea
$I = B \bmod 16$ \cite{stallings2019,wikipedia-cpu-cache}, y ahí se
concentra el daño: las palabras 0, 16 y 32 (índice 0) y las 1 y 17
(índice 1) compiten por una sola línea. Cinco de los seis fallos de
conflicto del directo pertenecen al primer grupo (los accesos 5, 12,
17, 18 y 19 de la traza) y el restante al segundo (acceso 10): cada
lectura expulsa a la anterior, que volverá a fallar cuando se pida de
nuevo. Es el patrón de \emph{thrashing} que Stallings ilustra tras la
Fig.~5.8 de su cap.~5 \cite{stallings2019} y que la clasificación de 3C
asocia a los fallos de conflicto \cite{wikipedia-metrics}; el marco del
conjunto de trabajo formaliza por qué alternar entre conjuntos disjuntos
desborda la caché \cite{denning1968,denning1972}; es también el
punto en que la caché deja de compensar \cite{null2024}.

El por conjuntos contiene el daño porque el conjunto 0 tiene dos vías:
solo tres accesos (los 17--19) no encuentran sitio y el bit de LRU
decide cuál de las dos palabras sale \cite{wikipedia-cpu-cache,hennessy2019}.
La asociativa total no lo sufre: con 16 líneas libres, cualquier palabra
encuentra hueco \cite{stallings2019}.

El resultado depende de la traza, y conviene decirlo. Si esta
alternara de forma exclusiva las palabras 0, 16 y 32, el mapeo directo
no acertaría ningún acceso tras los obligatorios y el por conjuntos
tampoco (tres palabras para dos vías con LRU es la secuencia clásica
que agota el reemplazo), mientras que la asociativa total mantendría las
tres residentes y solo pagaría sus tres primeros fallos. La traza real
intercala el resto de direcciones y por eso las tasas se quedan en
35--65\,\% en lugar de colapsar.

\subsection{Latencia de acceso y tiempo medio}

Los VCD permiten medir el tiempo que va desde que el banco emite
\texttt{req} hasta que \texttt{ready} valida el dato: 30~ns (3 ciclos
de reloj de 10~ns) cuando hay acierto y 60~ns (6 ciclos) cuando hay
fallo, valores idénticos en los tres mapeos porque comparten la misma
máquina de estados. La penalización de fallo es, por tanto,
$P = 30$~ns. Con el ratio $h$ de cada mapeo, el tiempo medio de acceso
queda
\[
  \mathrm{AMAT} \;=\; T_a + (1-h)\cdot P \,,
\]
con $T_a$ el tiempo de acierto; es la misma estructura que el tiempo
medio $c + (1-h)\cdot m$ de Tanenbaum y Austin \cite{tanenbaum2013} y
que los modelos de rendimiento de Stallings \cite{stallings2019} y de
Null y Lobur \cite{null2024}.

\begin{table}[t]
\centering
\small
\begin{tabular}{@{}lrrrr@{}}
\hline
\noalign{\smallskip}
 & $h$ & $1-h$ & $(1-h)\cdot P$ (ns) & AMAT (ns) \\
\noalign{\smallskip}\hline\noalign{\smallskip}
Directo & 0,35 & 0,65 & 19,5 & 49,5 \\
Por conjuntos & 0,50 & 0,50 & 15,0 & 45,0 \\
Total & 0,65 & 0,35 & 10,5 & 40,5 \\
\noalign{\smallskip}\hline
\end{tabular}
\caption{Tiempo medio de acceso por mapeo con $T_a = 30$~ns y
penalización de fallo $P = 30$~ns, medidos sobre los VCD.}
\label{tab:amat}
\end{table}

Con $T_a = 30$~ns, la Tabla~\ref{tab:amat} muestra que la asociativa
total ahorra 9~ns por acceso respecto a la directa (40,5 frente a
49,5~ns: cada acceso tarda un 22\,\% más en el mapeo directo) y que el
por conjuntos se queda a mitad de camino. Las duraciones de las tres
simulaciones confirman la cuenta: la parte fija del banco (arranque y
huecos entre accesos) es de 425~ns en los tres casos, de modo que
$425 + 20 \times \mathrm{AMAT}$ reproduce exactamente los 1415, 1325 y
1235~ns medidos en la Tabla~\ref{tab:resultados}.

Conviene matizar que $T_a$ es igual para los tres por decisión de
diseño: se comparan implementaciones funcionales, no silicio
(sección~\ref{sec:propuesta}). En hardware real, resolver un acierto
en la asociativa total exige comparar las 16 etiquetas en paralelo, con
más área y más energía que el único comparador del mapeo directo
\cite{hennessy2019,wikipedia-cpu-cache}; y como el AMAT tiende a $T_a$
cuando $h$ se acerca a 1, a partir de cierto punto añadir
asociatividad deja de rentabilizarse \cite{hennessy2019}.

Como ejercicio de sensibilidad, si la penalización de fallo fuera de
100~ns en lugar de 30~ns, el AMAT sería de 95, 80 y 65~ns: el mapeo
directo tardaría un 46\,\% más por acceso que el total, frente al
22\,\% del modelo de un ciclo. El orden no cambia, pero sí el valor en
juego: cuanto más cara es la memoria, más pesa la tasa de fallos y más
se justifica invertir en asociatividad \cite{hennessy2019}.

\subsection{Del tipo de fallo al tiempo de simulación}

\begin{table}[t]
\centering
\small
\begin{tabular}{@{}lrrrr@{}}
\hline
\noalign{\smallskip}
 & Todo acierto & Obligatorios & Conflicto & Total \\
 & (ns) & (+ns) & (+ns) & (ns) \\
\noalign{\smallskip}\hline\noalign{\smallskip}
Directo & 1025 & 210 & 180 & 1415 \\
Por conjuntos & 1025 & 210 & 90 & 1325 \\
Total & 1025 & 210 & 0 & 1235 \\
\noalign{\smallskip}\hline
\end{tabular}
\caption{Tiempo de simulación desglosado por tipo de fallo (parte fija
de 425~ns más 600~ns si los 20 accesos acertaran todos).}
\label{tab:desglose}
\end{table}

La Tabla~\ref{tab:desglose} lleva la clasificación de la
sección~\ref{sub:3c} al tiempo: si los 20 accesos hubieran acertado,
cada simulación habría durado 1025~ns (425~ns de arranque y huecos del
banco, más 600~ns de 20 aciertos de 30~ns). A esa base, los 7 fallos
obligatorios suman 210~ns en los tres mapeos, por igual, y los de
conflicto añaden 180, 90 y 0~ns. Las columnas cuadran con los 1415,
1325 y 1235~ns medidos: todo lo que una simulación tarda de más que
otra procede de los fallos de conflicto que su reparto de direcciones
provoca.

\subsection{Alcance de la comparación}

Los números describen esta traza y estos parámetros, no un rendimiento
global: la traza se diseñó con localidad temporal y espacial y con un
patrón deliberado de conflictos (sección~\ref{sec:propuesta}), la caché
es de 64~B y la memoria de 1~KB, y no aparecen fallos de capacidad
porque el conjunto de trabajo es pequeño. Con una traza más larga y
variada, la componente de capacidad crecería y también penalizaría a la
asociativa total \cite{bryant2015,patterson2017}. A ello se suma que
todo el flujo es simulación funcional con GHDL, sin síntesis
\cite{ghdl-docs}: los tiempos de esta sección son de simulación y de
un modelo de memoria de un ciclo, no de un procesador real.

A esto se suma que los parámetros elegidos condicionan el resultado: el
bloque es de una sola palabra (4~B), de modo que cuatro palabras
contiguas ocupan cuatro líneas distintas y no se aprovecha la localidad
espacial dentro del bloque \cite{bryant2015,wikipedia-locality,lam1991}; el
tamaño de bloque es precisamente uno de los parámetros que más afecta al
rendimiento de una caché \cite{smith1987}. Con
16 líneas el conjunto de trabajo cabe entero. Con bloques mayores o con
menos líneas, el reparto de fallos sería distinto
\cite{bryant2015,patterson2017}.

% ============================================================
% 6. CONCLUSIONES  (1 cuartilla)
% ============================================================
\section{Conclusiones}
\label{sec:conclusiones}

\textbf{Logros frente a los objetivos.} Este trabajo se marcó dos
objetivos (sección~\ref{sec:intro}): comprobar por simulación el
comportamiento teórico de cada mapeo y cuantificar sus diferencias
sobre una carga común, y ambos se cumplen. Se implementaron en
VHDL-2008 la memoria principal y
los tres mapeos de caché con una interfaz común, y cinco bancos de
prueba con autochequeo verifican cada dato leído, cada resultado
acierto/falto y los totales de la traza. En condiciones idénticas
(misma traza, misma memoria, mismo arbitraje), la tasa de aciertos
resultante es de 35\,\% en el mapeo directo, 50\,\% en el por conjuntos
y 65\,\% en el asociativo total, con un tiempo medio de acceso de 49,5,
45,0 y 40,5~ns respectivamente: los 6, 3 y 0 fallos de conflicto son la
única diferencia entre ellos, ya que los 7 fallos obligatorios y los
fallos de capacidad son comunes.

\textbf{Cuándo conviene cada mapeo.} El mapeo directo sigue siendo la
opción cuando mandan la simplicidad, el coste y un tiempo de acierto
bajo y previsible, por ejemplo en cachés pequeñas o en los primeros
niveles de una jerarquía modesta \cite{stallings2019,null2024}. El
asociativo por conjuntos es el compromiso habitual: con solo dos vías y
un bit de LRU por conjunto recupera buena parte de los aciertos
perdidos --en nuestra traza, 3 de los 6 fallos de conflicto-- a un
coste muy superior al del directo \cite{hennessy2019,wikipedia-cpu-cache}.
El asociativo total es el que más aciertos obtiene y conviene cuando el
presupuesto de tiempo, área y energía lo permite, porque exige
comparar todas las etiquetas en paralelo \cite{hennessy2019}.

\textbf{Limitaciones.} La caché es de solo lectura: no hay bit de
suciedad ni política de escritura, de modo que no se observan
invalidaciones ni actualizaciones \cite{bryant2015,wikipedia-cache-computing}.
Los tamaños (64~B de caché frente a 1~KB de memoria) y la traza de
20 accesos son pequeños para observar fallos de capacidad, y la
memoria de comportamiento tiene una latencia de un ciclo. Por último,
GHDL se usa como simulador --el backend mcode del equipo no permite
síntesis \cite{ghdl-docs}--: los resultados son funcionales, no de
silicio.

\textbf{Trabajo futuro.} Los siguientes pasos naturales son: incorporar
el bit de suciedad con las políticas \emph{write-through} y
\emph{write-back} \cite{wikipedia-cache-computing,bryant2015,stallings2019};
añadir un segundo nivel de caché para estudiar una jerarquía real
\cite{tanenbaum2013}; sustituir LRU por aproximaciones cuando crezca la
asociatividad, como FIFO, aleatorio o pseudo-LRU \cite{hennessy2019};
y ejecutar trazas de programas reales, más largas y con más bloques
destinados a hacer aparecer los fallos de capacidad
\cite{patterson2017,bryant2015}. Ampliar también la comparación a más
configuraciones --distinto número de vías y de líneas-- para obtener la
curva de asociatividad frente a tasa de aciertos \cite{bryant2015}.

\textbf{Reproducibilidad.} Todo el experimento se repite con los
wrappers del repositorio: \texttt{scripts/\allowbreak vhdl-sim}
vuelve a generar los cinco resultados y
\texttt{scripts/\allowbreak latexbuild} recompone este informe, mientras
la integración continua ejecuta los cinco bancos en
cada \emph{pull request}. Los valores de este artículo no dependen,
por tanto, de una ejecución aislada, sino del propio código.

% ============================================================
% REFERENCIAS  (>= 15 — indicar exactamente dónde aparece el texto)
% ============================================================
% IMPORTANTE: cada referencia debe permitir ubicar exactamente dónde
% aparece el texto sugerido (capítulo/§/página, o URL + fecha de consulta).
% Anota la ubicación en el campo note de la entrada.
\bibliographystyle{splncs04}
\bibliography{referencias}

% ============================================================
% APÉNDICE — CÓDIGO VHDL
% ============================================================
\appendix
\section{Código VHDL}
\label{app:codigo}

Código completo de los cuatro módulos de la unidad: memoria principal
(Listado~\ref{lst:memoria}), mapeo directo (Listado~\ref{lst:directo}),
asociativo total (Listado~\ref{lst:asociativa}) y asociativo por
conjuntos (Listado~\ref{lst:conjuntos}). Los listados se incluyen con
\texttt{lstinputlisting} directamente desde
\texttt{vhdl/\allowbreak mapeo-\allowbreak memoria/\allowbreak src/},
de modo que siempre coinciden con el código que ejecutan los bancos de
prueba de la sección~\ref{sec:propuesta}.

\lstinputlisting[
  caption={memoria\_principal.vhd --- RAM de comportamiento de 1~KiB
    (256 palabras de 4~B) con patrón de fábrica y puerto de escritura},
  label={lst:memoria}]{../../vhdl/mapeo-memoria/src/memoria_principal.vhd}

\lstinputlisting[
  caption={cache\_directo.vhd --- caché de mapeo directo: 16 líneas de
    4~B, sin reemplazo, FSM de cuatro estados},
  label={lst:directo}]{../../vhdl/mapeo-memoria/src/cache_directo.vhd}

\lstinputlisting[
  caption={cache\_asociativa.vhd --- caché asociativa total: 16 líneas
    con comparación paralela de las etiquetas y LRU por contador},
  label={lst:asociativa}]{../../vhdl/mapeo-memoria/src/cache_asociativa.vhd}

\lstinputlisting[
  caption={cache\_conjuntos.vhd --- caché por conjuntos: 8 conjuntos de
    2 vías con LRU de 1 bit por conjunto},
  label={lst:conjuntos}]{../../vhdl/mapeo-memoria/src/cache_conjuntos.vhd}

\end{document}
